\section{Related Work}

Below, we review work on enhancing the accessibility of social VR environments and supporting BLV individuals with guides and visual interpretation in the physical world.

\subsection{Enhancing the Accessibility of Social VR Environments}

Social virtual environments are typically complex, requiring users to navigate surroundings, interpret avatar movements, and discern nonverbal cues. Researchers have explored many approaches to making these environments accessible, often by adding sonification or haptic effects to objects and users’ actions \cite{ji2022vrbubble, balasubramanian2023sceneweaver, jain2015hmd, goncalves2023navaware, killough2025vrsight, killough2025vrsightdemo}. For instance, systems like Ji et al.’s VRBubble \cite{ji2022vrbubble} and those by Killough et al. \cite{killough2025vrsight, killough2025vrsightdemo} use spatial audio feedback to convey the number, proximity, and location of avatars and objects in VR scenes. Others have focused on translating nonverbal cues, like eye contact and head nodding, into audio signals and haptic vibrations \cite{wieland2022nvc, wieland2023gaze, segal2024socialcues, dang2023music, collins2023gaze, jung2024accessible}. While crucial for low-level perception, these discrete sensory solutions often become sensorially overwhelming in the dense landscapes of social VR, and they fail to provide the high-level information necessary for task completion and social fluency. Our work seeks to address these issues by providing a guide tailored to provide such high-level output for complex social environments.

Enhancing the accessibility of others’ location and behavior addresses some aspects of VR experiences, but it is equally important to consider how disabled users wish to represent themselves during these experiences. To this end, researchers have explored how disabled people wish to present in social virtual worlds \cite{gualano2023invisible, gualano2024looking, zhang2023signifier, zhang2022selfpres, mack2023avatarrep}. Zhang et al. highlighted the lack of disability representation in VR avatar customization, with some disabilities being excluded from menus and others given only one option, such as a hearing aid for d/Deaf people \cite{zhang2022selfpres}. Gualano et al. explored how people with invisible disabilities like neurodivergence or chronic pain wish to represent their disabilities, finding that symbolic representations like social batteries were desired \cite{gualano2023invisible}. Our work complements this stream of research by exploring how an AI guide can provide assistance while offering new forms of disability representation through its own embodiments.

Among efforts to enhance social VR accessibility, our research focuses on exploring virtual guides to support BLV users. In our prior work by Collins and Jung et al., we introduced a framework featuring a remote human assistant who served as both a conversational partner and a tool for navigation and visual interpretation \cite{collins2023guide}. In this study, BLV users explored virtual parks with the assistance of this human guide. This human guide framework demonstrated high effectiveness in helping BLV users explore virtual parks. However, this approach is fundamentally limited by the availability of human assistants, users’ wish to be independent, and difficulty in matching guiding styles to individual preferences.

To address the limitations of human guides, our subsequent work by Collins et al. described the design of an AI guide with six combinations of speech characteristics and appearances called “personas.” \cite{collins2024ai} These personas offer varying levels of social interaction and fantastical embodiment, including a friendly guide dog and blunt white cane, catering to diverse user needs and providing new avenues for disability representation in social VR. However, this work also has several limitations. First, this guide was not tested with BLV users in social VR, limiting our understanding of how BLV people utilize it. Second, while the guide was co-designed with BLV members of our research team, it lacked broader design input from a general BLV audience. Finally, we did not provide explicit details about the LLM prompt used to enable our guide, which is crucial to others’ ability to create similar AI-powered assistants for BLV users.

Thus, we improved the initial design of our AI guide prototype, focusing on three key personas. We incorporated key modifications based on new pilot studies and an updated design process, and shared our full AI prompt as supplemental material. Our current work moves beyond our prior design focus by providing the first empirical evidence of how BLV individuals utilize and interact with an AI guide in social VR. By incorporating their feedback, we identified opportunities to build reliable, user-centered guides informed by the unique ways BLV individuals interact with AI assistance in VR.

\subsection{Supporting BLV People with Guides in the Physical World}

Guides are a widespread tool used to support BLV people in the physical world. Researchers have examined the use of guides by BLV people for decades \cite{kayukawa2020blindpilot, garaj2003pedestrian, guerreiro2019cabot, lannan2019campus, jagt2023familiar}. A “guide” can be any assistant providing physical support to navigate from one area to another. This includes dog guides (commonly referred to as “guide dogs”), robots, and sighted human guides (assistants offering physical navigation assistance). 

Sighted human guides and guide dogs have been widely studied for their dual roles as companions and navigators \cite{ball2022runners, miner2001dogguide, vincenzi2021interdependence, small2015interconnecting, macpherson2016guiding}. For instance, Small found via autoethnographic entries and questionnaires that BLV tourists fostered companionship with sighted tourists who guided them, allowing them to enjoy visual aspects of tourism, like artwork, through informal socializing \cite{small2015interconnecting}. Miner et al. conducted open-ended interviews with eight BLV people about working with guide dogs, finding that guide dogs increased their handlers’ social confidence and acted as icebreakers in social spaces as well as met navigational needs \cite{miner2001dogguide}.

Although the companionship of sighted guides and guide dogs is valued, less-social guidance robots have also emerged as a topic of interest among HCI researchers \cite{kayukawa2020blindpilot, guerreiro2019cabot, azenkot2016servicerobots, kayukawa2023museum, bonani2018robots, hwang2024lessons, hwang2024towards}. For example, Hwang et al. found that with the increased availability of robots, robot guides could serve as less expensive alternatives to guide dogs if they borrow best practices from guide dogs \cite{hwang2024lessons}. In follow-up work, Hwang et al. conducted interviews and observed guidance sessions with five guide dog trainers and 23 guide dog handlers to identify these practices to inform future robot guides \cite{hwang2024towards}.

Given the proven utility of guides in complex situations, we aimed to develop virtual guides for VR by incorporating established best practices from their physical-world counterparts.

\subsection{Supporting BLV People with Visual Interpretation in the Physical World}

Another key tool to support BLV people in the physical world are the well-studied visual interpretation services \cite{gonzalez2024scenedesc, gonzalez2025towards, garaj2003pedestrian, lannan2019campus, bigham2010vizwiz}. These services allow BLV people to capture images or video of their surroundings, often via camera-based phone applications, and receive verbal descriptions of the captured media from human or AI interpreters.

One well-known visual interpretation application is VizWiz, an application that connects BLV users to remote sighted workers who answer queries about provided images \cite{bigham2010vizwiz}. Similar human-powered visual interpretation services that have been studied include Aira and BeMyEyes \cite{garaj2003pedestrian, lannan2019campus}. For instance, Lannan interviewed nine BLV college students about how they used Aira, finding that the application helped them engage in group activities as well as navigate their campuses \cite{lannan2019campus}.

In recent years, some researchers have begun focusing on the use of AI-powered visual interpretation, driven by the increased availability and quality of descriptions generated by LLMs \cite{gonzalez2024scenedesc, gonzalez2025towards, lin2025viptour}. For example, Gonzalez and Collins et al. conducted a two-week diary study of an LLM-powered interpretation application, followed by interviews with 16 BLV people. They found BLV people preferred AI in certain use cases over human interpretation. These included reading financial documents and taking pictures in private areas like bathrooms since AI posed fewer privacy concerns \cite{gonzalez2024scenedesc}. Lin et al. explored using an AI-driven system to help 30 BLV participants navigate unfamiliar environments, finding that BLV people prioritized obstacle avoidance, wayfinding, and understanding environmental aesthetics distilled in a concise, personalized manner \cite{lin2025viptour}.

Given the many uses of visual interpretation tools, we aimed to create an AI guide that is able to respond to all types of queries, and can provide reliable enough information that users feel comfortable utilizing it even for private scenarios.